% ==========================================================================
%  Chapter 10 — Conclusion and Future Work
% ==========================================================================
\chapter{Conclusion and Future Work}
\label{ch:conclusion}

\section{Summary of Contributions}
\label{sec:conclusion-summary}

This report has presented the v2 redesign of User-Mode Linux, the most
significant rearchitecture of the \UML{} subsystem since Jeff Dike's
original implementation in 2000~\cite{dike2000}.  The contributions
span architecture, implementation, validation, tooling, and
documentation:

\begin{enumerate}[leftmargin=2em]
  \item \textbf{Backend operations abstraction.}
    \texttt{struct um\_backend\_ops} decouples the \UML{} kernel from
    any single syscall-interposition mechanism, enabling pluggable
    backends with a tested behavioral contract.

  \item \textbf{The \KVMvTwo{} backend.}
    A complete \KVM{}-based execution backend that replaces the
    ptrace/seccomp signal-mediated syscall path with
    \texttt{VMRESUME}/\texttt{VMEXIT}.  The LSTAR-gadget fast path
    achieves $\sim$90-cycle syscall latency, a $3$--$4\times$ improvement
    over seccomp on real workloads and up to $900\times$ on
    microbenchmarks.

  \item \textbf{SMP support.}
    Multi-processor execution with per-vCPU dispatch, IPI emulation,
    TLB shootdown, deferred-free TLB batching, and memory barrier
    semantics.

  \item \textbf{Snapshot and record/replay.}
    KVM-aware snapshot (\SI{38.9}{\micro\second} capture,
    \SI{13}{\micro\second} restore) and a record/replay state machine
    with static-key gating, enabling deterministic debugging and
    fast-restart fuzzing workflows.

  \item \textbf{Composable instrumentation.}
    KASAN, KMSAN, KCSAN, KFENCE, KCOV, kprobes, ftrace, and BPF JIT,
    all selectable through nine named profiles from a single tree.
    Runtime-flippable hooks via static keys allow post-failure
    observability without rebuilding.

  \item \textbf{Comprehensive validation.}
    53 kselftest suites, the \texttt{umlctl mission} six-phase
    acceptance gate, CPython substrate parity (21/21 on both backends),
    and cross-host benchmarking on eight machines.

  \item \textbf{\texttt{umlctl}: a Rust launcher.}
    A 22,000-line Rust CLI that reduces the barrier from a page of
    shell commands to a single invocation, with declarative deployment,
    cgroup v2 management, and reproducible acceptance testing.

  \item \textbf{Documentation.}
    134,500 lines of operator documentation, redesign planning
    materials, and upstream submission notes.
\end{enumerate}

\noindent
The seven series-level commits on \texttt{torvalds/master} (v7.1-rc5+)
represent a cohesive body of work: each series is independently
reviewable and bisectable, and the complete set brings \UML{} from a
largely unmaintained architecture port to a viable platform for kernel
research, debugging, fuzzing, and sandboxed execution.


\section{Current State (May 2026)}
\label{sec:conclusion-current}

As of the reporting cutoff, the following capabilities are operational:

\begin{description}[style=nextline,leftmargin=2em]
  \item[Both backends pass the full CPython test suite.]
    The substrate parity gate produces identical results on \KVMvTwo{}
    and seccomp: PASS=25, FAIL=3, XFAIL=3.  The three failures are
    known \UML{} limitations (not regressions), and the three expected
    failures are test-side issues unrelated to the backend.

  \item[\KVMvTwo{} is 3--4$\times$ faster on real workloads.]
    The Python startup benchmark shows 2.99--4.06$\times$ speedup over
    seccomp across eight host types.  The SMP stress tier shows
    3.71--8.72$\times$ speedup.  Microbenchmarks show 193--908$\times$
    speedup, reflecting the elimination of the host-exit round-trip for
    gadget-covered syscalls.

  \item[SMP works.]
    The mt-mini stress test runs with T=8 threads on ncpus=4 virtual
    CPUs with strict memory verification and zero failures.

  \item[Snapshot and record/replay are implemented.]
    The \KVM{}-aware snapshot surface captures in \SI{38.9}{\micro\second}
    and restores in \SI{13}{\micro\second} median.  The record/replay
    state machine logs SYSCALL, RDTSC, and SIGALRM events with
    static-key gating.  KUnit suites cover 3/3 snapshot cases and
    7/7 record/replay cases.

  \item[\texttt{umlctl mission} passes 6/6 phases.]
    The canonical acceptance gate completes in $\sim$15 minutes with
    100/100 soak iterations passing, zero panics, and a JSON
    scoreboard for audit.
\end{description}


\section{What Remains}
\label{sec:conclusion-remaining}

Despite the breadth of the current implementation, several significant
work items remain before the redesign can be considered complete.

\subsection{Phase J: 24-Hour Continuous Soak}

The \texttt{umlctl mission} gate provides 15-minute validation with
100 soak iterations.  The 24-hour long-tail soak---running thousands
of iterations across all workload templates---is the remaining
wall-clock evidence needed for milestone candidates.  The rig is in
place; the remaining cost is compute time.

\subsection{Upstream Acceptance}

The seven series are drafted and validated, but none has been
submitted to LKML.  The submission strategy
(Section~\ref{sec:v2-upstream}) is staged: small independent patches
first, the tentpole RFC second, and the \KVM{} backend last.  The
primary risk is review bandwidth: the backend-ops abstraction touches
core \UML{} code and requires buy-in from the \UML{} maintainer and
the broader \texttt{linux-arch} community.

\subsection{Architecture Ports}

The v2 redesign targets x86\_64 exclusively.  Two additional
architecture ports are planned:

\begin{description}[style=nextline,leftmargin=2em]
  \item[ARM64 (\texttt{aarch64}).]
    ARM64 \KVM{} support is mature, and the backend-ops abstraction is
    architecture-neutral by design.  The primary work is implementing
    the ARM64-specific portions of the \KVMvTwo{} backend: exception
    vector setup, register marshaling, and the equivalent of the
    LSTAR gadget (using \texttt{HVC} or \texttt{SMC} instructions for
    fast-path trapping).

  \item[RISC-V.]
    RISC-V \KVM{} support is upstream and stabilizing.  The port
    follows the same pattern as ARM64, with RISC-V-specific trap
    handling via the \texttt{ecall} mechanism.
\end{description}

\subsection{syzkaller Integration}

The in-tree forkserver surface is pinned (snapshot v1, contract
C-09).  The remaining work is the \texttt{pkg/vm/uml} Go driver in
the syzkaller repository~\cite{syzkaller}, which maps syzkaller's VM
management API to \UML{} instance lifecycle.  This is an external
contribution, not a kernel-tree change.

\subsection{Record/Replay Host-Side Wiring}

The record/replay log shape, observe/consume API, and KUnit suite are
proven.  The remaining work is the host-side signal-interception path
that converts a real \texttt{SIGALRM} into an
\texttt{observe\_sigalrm} call, plus the trap-on-RDTSC VMCS
programming for deterministic timestamp counters.

\subsection{Additional Validation}

\begin{itemize}[leftmargin=2em]
  \item \textbf{LTP curation.}  A curated KEEP/SKIP list and
    \UML{}-appropriate runner refinements for the Linux Test Project.
  \item \textbf{CPUID feature un-masking.}  AVX-512, AMX, PKRU, and
    CET un-masking following the SMP-T77 checklist.
  \item \textbf{Snapshot ELF64-core export.}  Lifting snapshot data
    to a gdb/crash-tool-loadable artifact.
\end{itemize}


\section{The Broader Impact}
\label{sec:conclusion-impact}

\UML{} occupies a unique position in the virtualization landscape---one
that no other technology can fill.

\subsection{Full Kernel, Full Debugging, No Hardware}

Every competing approach sacrifices at least one of these three
properties:

\begin{itemize}[leftmargin=2em]
  \item \KVM{}/QEMU provides a full kernel and hardware acceleration,
    but debugging requires GDB stubs, serial consoles, or JTAG
    probes.
  \item gVisor provides sandboxing without hardware requirements, but
    it is not the Linux kernel---semantic divergences are inherent.
  \item Containers provide the real kernel with low overhead, but they
    share the host kernel and cannot isolate kernel bugs.
  \item \textsc{LKL} provides the real kernel code at library-call
    speed, but sacrifices process isolation and scheduling fidelity.
  \item Unikernels sacrifice multi-process semantics entirely.
\end{itemize}

\noindent
\UML{} is the only technology that provides all three: the complete
Linux kernel, running with full \texttt{gdb} debuggability, without
requiring hardware virtualization extensions.  When the host does
provide \texttt{/dev/kvm}, the \KVMvTwo{} backend makes \UML{}
competitive in performance while retaining the debugging story.

\subsection{The Cloud-Native and eBPF Context}

The rise of eBPF~\cite{ebpf} as a kernel programmability mechanism
has created a new class of kernel developers who write BPF programs
but rarely compile or debug the kernel itself.  \UML{} can serve as
the bridge: a BPF developer can run a \UML{} instance, load their
programs, observe the kernel's internal state with kprobes and
ftrace, and diagnose issues that are invisible from the eBPF
toolchain alone.

Similarly, the cloud-native ecosystem increasingly relies on kernel
features (cgroups, namespaces, seccomp, io\_uring) without direct
access to the kernel development workflow.  \UML{} provides that
access in a form factor that fits a CI/CD pipeline: compile, boot,
test, and tear down in seconds, without a VM, without hardware
access, without root privileges.

\subsection{Security and Reproducibility}

The security research community needs two things that are difficult
to combine: a realistic attack surface and a controlled execution
environment.  \UML{} provides both.  The attack surface is real---it
is the Linux kernel's syscall interface---and the environment is
controlled: deterministic time, record/replay, and snapshot/restore
allow a researcher to reproduce and analyze an exploit as many times
as needed.  No serial console, no JTAG probe, no hardware testbed.


\section{Closing Thoughts}
\label{sec:conclusion-closing}

\UML{} was created 26 years ago to solve a simple problem: make the
Linux kernel debuggable.  For a time, it was the only way to run a
kernel under \texttt{gdb}.  Then \KVM{} arrived, then containers, then
gVisor, then eBPF, and \UML{} faded into maintenance mode.  The
underlying problem---making the kernel observable, testable, and
debuggable without dedicated hardware---never went away.  It just
acquired more competitors, none of which fully solve it.

The window for \UML{}'s revival is open now, and it is not large.

The Linux Kernel Library~\cite{lkl} demonstrates that a
library-compiled kernel can achieve extraordinary fuzzing throughput.
Rust VMM crates (rust-vmm, crosvm~\cite{crosvm}, Cloud
Hypervisor~\cite{cloud-hypervisor}) are making it trivial to build
custom hypervisors.  eBPF~\cite{ebpf} is absorbing an increasing
share of the kernel observability use case.  If \UML{} does not
catch up---if it remains a curiosity with a ptrace-era performance
profile and no modern tooling---these technologies will collectively
eat the niche.

The v2 redesign is an attempt to close the gap before it becomes
permanent.  The backend-ops abstraction makes \UML{} extensible.
The \KVMvTwo{} backend makes it fast.  The profile system makes it
composable.  The snapshot and record/replay infrastructure make it
useful for fuzzing and debugging.  The \texttt{umlctl} launcher
makes it accessible.

Whether this attempt succeeds depends on upstream acceptance.  The
code is written, the tests pass, the benchmarks are favorable, and
the documentation is comprehensive.  What remains is the social
process of kernel review: demonstrating to the \UML{} maintainer
and the broader kernel community that this work is worth maintaining,
that the backend-ops contract is the right abstraction, and that
the \KVMvTwo{} backend justifies its complexity.

The answer to ``why \UML{}?'' has not changed in 26 years: because
the best way to understand a kernel is to run it where you can poke
at it.  The v2 redesign simply ensures that the poking is fast enough,
instrumented enough, and accessible enough to be worth doing in 2026.
